\documentclass[11pt]{article}

% Pacotes
\RequirePackage{luatex85}
\usepackage{amsmath}
\usepackage{amssymb}
\usepackage{graphicx}
\usepackage[colorlinks=true, allcolors=blue]{hyperref}
\usepackage{amsthm}
\usepackage{framed, color,xcolor}
\usepackage{stmaryrd}
\usepackage{rotating}           % Usado para rotacionar o texto
\usepackage[all,knot,arc,import,poly]{xy}   % Pacote para desenhos gráficos
% Este pacote pode conflitar com outros pacotes gráficos como o ``pictex''
% Então é necessário usar apenas um dos pacotes conflitantes
\usepackage{tikz, tikz-cd}
\usepackage{epigraph}
\usetikzlibrary{babel}
\usetikzlibrary{positioning}
\usetikzlibrary{calc}
\usepackage[stix2]{newtxmath}
\usepackage{mathtools, nccmath}
%\usepackage[usenames,dvipsnames]{xcolor}
\usepackage{amssymb, mathrsfs}
\usepackage[numbers]{natbib}
\usepackage{fontspec}
\usepackage{enumitem} % pra substituir o símbolo do itemize por alguma coisa legal, tipo setinhas:  \begin{itemize}[label=\ding{212}]
\usepackage{pifont}
\usepackage{stmaryrd} % for \arrownot (notação de adjacência em grafo)
\usepackage{quiver} %pacote do quiver pra desenhar diagramas mais legais

%Fontes e Línguas
\usepackage[portuguese]{babel}
%\usepackage{eufrak}
\usepackage{emoji}
\setmainfont{DarwinSerif-Regular.otf}
\newfontfamily\fakebold[FakeBold=1.75]{DarwinSerif-Regular.otf}
\newfontfamily\fakeitalic[FakeSlant=0.2]{DarwinSerif-Regular.otf}
\renewcommand{\textbf}[1]{{\fakebold#1}}
\renewcommand{\textit}[1]{{\fakeitalic#1}}
\renewcommand{\bfseries}{\fakebold}
\renewcommand{\itshape}{\fakeitalic}



%Cores


%Comandinhos que gosto
\usepackage{mathtools}
\newcommand{\defeq}{\vcentcolon=}
\newcommand{\eqdef}{=\vcentcolon}
\newcommand{\bigslant}[2]{{\raisebox{.2em}{$#1$}\left/\raisebox{-.2em}{$#2$}\right.}}
\newcommand{\logo}{\Rightarrow}
\newcommand{\pulinho}{\vspace{.5cm}}
\newcommand{\edge}{%
  \mathrel{-}% minus as a relation
  \joinrel\joinrel % some backup
  \mathrel{-}% minus as a relation
}
\newcommand{\notedge}{%
  \mathrel{\mkern2mu}% a small advancement
  \arrownot % a short slash
  \mathrel{\mkern-2mu}% compensate
  \edge
}

\newcommand{\Aut}{\text{Aut}}
\newcommand{\Cay}{\text{Cay}}

\newcommand{\luafacil}{
\emoji{waxing-crescent-moon}
}

\newcommand{\luamedio}{
\emoji{first-quarter-moon}
}

\newcommand{\luadificil}{
\emoji{waxing-gibbous-moon}
}

\newcommand{\luaprojeto}{
\emoji{full-moon}
}

\let\varphi\phi

%Caixinhas coloridas

\theoremstyle{definition}
\newtheorem{theorem}{Teorema}
\newenvironment{teorema}{
\definecolor{shadecolor}{HTML}{f59aa3}
\begin{shaded}
\begin{theorem}
}
{
\end{theorem}
\end{shaded}
}

\theoremstyle{definition}
\newtheorem*{theorem*}{Teorema}
\newenvironment{teorema*}{
\definecolor{shadecolor}{HTML}{f59aa3}
\begin{shaded}
\begin{theorem*}
}
{
\end{theorem*}
\end{shaded}
}

\theoremstyle{definition}
\newtheorem{prop}{Proposição}
\newenvironment{proposicao}{
\definecolor{shadecolor}{HTML}{FFDBDF}
\begin{shaded}
\begin{prop}
}
{
\end{prop}
\end{shaded}
}

\theoremstyle{definition}
\newtheorem*{prop*}{Proposição}
\newenvironment{proposicao*}{
\definecolor{shadecolor}{HTML}{FFDBDF}
\begin{shaded}
\begin{prop*}
}
{
\end{prop*}
\end{shaded}
}

\theoremstyle{definition}
\newtheorem{cor}{Corolário}
\newenvironment{corolario}{
\definecolor{shadecolor}{HTML}{f59aa3}
\begin{shaded}
\begin{cor}
}
{
\end{cor}
\end{shaded}
}

\theoremstyle{definition}
\newtheorem*{cor*}{Corolário}
\newenvironment{corolario*}{
\definecolor{shadecolor}{HTML}{f59aa3}
\begin{shaded}
\begin{cor*}
}
{
\end{cor*}
\end{shaded}
}



\theoremstyle{definition}
\newtheorem{definition}{Definição}
\newenvironment{definicao}{
\definecolor{shadecolor}{HTML}{defcfc}
\begin{shaded}
\begin{definition}
}
{
\end{definition}
\end{shaded}
}

\theoremstyle{definition}
\newtheorem*{definition*}{Definição}
\newenvironment{definicao*}{
\definecolor{shadecolor}{HTML}{defcfc}
\begin{shaded}
\begin{definition*}
}
{
\end{definition*}
\end{shaded}
}

\theoremstyle{definition}
\newtheorem{example}{Exemplo}
\newenvironment{exemplo}{
\definecolor{shadecolor}{HTML}{b9e6d3}
\begin{shaded}
\begin{example}
}
{
\end{example}
\end{shaded}
}

\theoremstyle{definition}
\newtheorem*{example*}{Exemplo}
\newenvironment{exemplo*}{
\definecolor{shadecolor}{HTML}{b9e6d3}
\begin{shaded}
\begin{example*}
}
{
\end{example*}
\end{shaded}
}

\theoremstyle{definition}
\newtheorem{veronica}{Lema}
\newenvironment{lema}{
\definecolor{shadecolor}{HTML}{FFF1F2}
\begin{shaded}
\begin{veronica}
}
{
\end{veronica}
\end{shaded}
}

\theoremstyle{definition}
\newtheorem*{veronica*}{Lema}
\newenvironment{lema*}{
\definecolor{shadecolor}{HTML}{FFF1F2}
\begin{shaded}
\begin{veronica*}
}
{
\end{veronica*}
\end{shaded}
}

\newenvironment{prova}{\paragraph{Demonstração:}}{\hfill$\square$ 
\pulinho
}

\newenvironment{aviso}{
\definecolor{shadecolor}{HTML}{fffe9a}
\begin{shaded}
}{\end{shaded}}

\theoremstyle{definition}
\newtheorem{exersisio}{Exercício}
\newenvironment{exercicio}{
\definecolor{shadecolor}{HTML}{CFB5FF}
\begin{shaded}
\begin{exersisio}
}
{
\end{exersisio}
\end{shaded}
}

%Tamanho do papel
\usepackage[a4paper,top=1cm,bottom=1cm,margin=3.5cm]{geometry}



%%%%%%%%%%%%%%%%%%%%%%%%%%%%%%%%%%%%%%%%%%%%%%%


%Definições só pra esse texto
\usepackage[bb=boondox]{mathalpha}



\title{Emaranhados}
\author{Violeta Mar}
\date{}


\begin{document}

\maketitle

A teoria de emaranhados (\textit{tangles}, em inglês) nasce na teoria de conectividade de grafos finitos. Aqui apresentamos uma versão `conjuntista' dela, vinda da versão simplificada apresentada por Diestel em seu livro cujo objetivo era aplicações em inteligência artificial. \cite{diestel2024tangles}

\begin{definicao*}[Emaranhados e consistência]Seja $V$ um conjunto, e $S \subset \mathcal{P}(V)$ um subconjunto fechado para complementar. Uma função $\tau: S \rightarrow \{0,1\}$ é dita um \textit{emaranhado} se satisfaz
  \begin{itemize}[label=\ding{212}]
  \item $\tau(A) = 1$ ou $\tau(\lnot A) = 1$ para qualquer $A \in S$
  \item \textit{consistência}: dados quaisquer $A,B,C \in S$, se a interseção destes três é vazia, então pelo menos um dos $\tau(A),\tau(B), \tau(C)$ é $0$.
  \end{itemize} 
\end{definicao*}

Vamos dizer que $\tau$ escolhe um $A$ se $\tau(A) = 1$. Veja que consistência implica que $\tau$ sempre escolhe apenas um dentre $A$ e seu complementar $\lnot A$.

Esta definição pode lembrar a definição de ultrafiltros para álgebras de Boole. De fato, se $S$ é uma álgebra de Boole, emaranhados são apenas ultrafiltros.

\begin{proposicao}[Emaranhados são ultrafiltros quando $S$ é de Boole]
  Suponha que $S$ seja fechado para união e interseção, ou seja, $S$ forma uma álgebra de Boole (ou corpo de conjuntos). Então $\tau:S \rightarrow \{0,1\}$ é um emaranhado se e somente se é um homomorfismo de álgebras de Boole, ou seja, um ultrafiltro.
\end{proposicao}
\begin{prova}
  Se $\tau$ é homomorfismo, então $\tau(A\cap B \cap C) = 0$ quando $A \cap B \cap C = \varnothing$, logo $\tau(A)\tau(B)\tau(C) = 0$, que só pode acontecer se pelo menos um destes valores é $0$.

  Reciprocamente, dado um emaranhado $\tau$ e dois elementos $A,B \in S$, considere os quatro casos de valores para $\tau(A)$ e $\tau(B)$. Faremos o primeiro: suponha $\tau(A) = \tau(B) = 1$. Como a interseção de $A$, $B$ e $\lnot (A \cap B)$ é vazia, devemos ter $\tau (\lnot (A \cap B)) = 0$ e portanto $\tau(A \cap B) = 1 = \tau(A)\tau(B)$. Como a interseção de $A \cap B$ e $\lnot A \cap \lnot B$ é vazia, então $\tau(\lnot A \cap \lnot B) = 0 $. Mas $\lnot (\lnot A \cap \lnot B) = A \cup B$ e logo $\tau(A \cup B) = 1 = \tau(A) + \tau(B)$. Fazemos contas análogas para cada um dos outros três casos e concluimos que $\tau$ é homomorfismo.
\end{prova}

Vamos deixar essa prova um pouco mais clara introduzindo alguns conceitos a mais.

\begin{definicao*}[Cantos e comparabilidade]
  Dados dois conjuntos $A,B \in S$, chamamos os quatro conjuntos $A \cap B, A \cap \lnot B, \lnot A \cap B, \lnot A \cap \lnot B$ de cantos do par $A,B$. Dizemos que $A$ e $B$ são comparáveis se pelo menos um destes conjuntos é vazio, caso contrário dizemos que são incomparáveis.
\end{definicao*}

Veja que $A$ é comparável a $B$ exatamente quando pelo menos uma das inclusões $A \subset B, A \subset \lnot B, \lnot A \subset B, \lnot A \subset \lnot B$ é verdadeira. Ou seja, $A$ e $B$ são incomparáveis quando $\{A, B, \lnot A, \lnot B\}$ é um conjunto de elementos que são incomparáveis em relação a ordem dada pela inclusão.

O valor de $\tau$ dos quatro cantos é inteiramente determinado pelo valor de $\tau$ em $A$ e em $B$ - é isto que essencialmente mostramos nos quatro casos da demonstração da Proposição $1$.

Os cantos de um par de conjuntos em $S$ podem não estar presentes em $S$. Dado um $S \subset \mathcal{P}(V)$ fechado pra complementar, dizer que todos os cantos de $S$ sempre estão em $S$ é equivalente a dizer que $S$ é uma álgebra de Boole. Na teoria de emaranhados, uma versão mais fraca vai ser suficiente para estabelecermos os resultados legais que buscamos.

\begin{definicao*}[Álgebras sub-Boole]
  Um subconjunto $S \subset \mathcal{P}(V)$ é dito sub-Boole (ou, na terminologia do Diestel, um conjunto submodular) se é fechado por complementar e dados quaisquer dois $A,B \in S$ temos que $A \cup B \in S$ ou $A \cap B \in S$.
\end{definicao*}

Temos dois pares de cantos que dizemos \textit{opostos}: $A \cap B$ e $\lnot A \cap \lnot B$; $\lnot A \cap B$ e $A \cap \lnot B$. Dado um $S \subset \mathcal{P}(V)$ fechado pra complementar, dizer que dentre dois cantos opostos, pelo menos um está em $S$ é equivalente a dizer que $S$ é sub-Boole.

Adicionando cantos, então, sempre conseguimos estender um dado $S$ para um $\tilde{S}$ que seja um sub-Boole ou uma de fato álgebra de Boole. Porém, um emaranhado $\tau$ de $S$ não necessariamente se estende para um emaranhado $\tilde{\tau}$ em $\tilde{S}$: veja, na definição, pedimos apenas que triplas de conjuntos com interseção vazia devem ter pelo menos um de seus conjuntos não escolhidos por $\tau$. Ao adicionarmos cantos ao $S$, não só triplas mas agora interseções de possivelmente até $6$ conjuntos devem ser consistentes para que tenhamos um emaranhado. Nem todo emaranhado vai ter esse nível mais alto de consistência. Por exemplo...

\begin{exemplo*}
  Seja $V = \{1,2,3,4\}$, $A = \{2,3,4\}, B = \{1,3,4\}, C = \{1,2,4\}, D = \{1,2,3\}$. Faça $S = \{A,B,C,D,\lnot A, \lnot B, \lnot C, \lnot D\}$. Então, a função $\tau$ que atribui $1$ a todos os $A,B,C,D$ e $0$ a seus complementares é um emaranhado. Mas, se adicionarmos apenas o canto $A \cap B$ a $S$, já não podemos estender $\tau$ para um $\tilde{\tau}$. Pois, se pudessemos, o valor de $\tilde{\tau}$ em $A \cap B$ deveria ser $1$, porém $(A \cap B) \cap C \cap D$ é vazio, e nenhum dos três $\tilde{\tau}(A \cap B), \tilde{\tau}(C), \tilde{\tau}(D)$ é $0$.
\end{exemplo*}

É claro, como o valor de um emaranhado $\tau$ em qualquer um dos quatro cantos de $A,B$ é determinado pelos valores de $\tau$ em $A$ e $B$, se existir uma extensão $\tilde{\tau}$ para $\tilde{S}$ então ela deve ser única.




\bibliographystyle{plain}
\bibliography{references}


\end{document}
